\section*{Bab \ref{ch:g8:the-atmosphere} --- Atmosfer}

\begin{solution}{exo:g8:the-atmosphere:1}
(a) \omterm{def:g8:the-atmosphere:layers}{stratosfer}; (b) \omterm{def:g8:the-atmosphere:layers}{troposfer}; (c) \omterm{def:g8:the-atmosphere:layers}{mesosfer}; (d) \omterm{def:g8:the-atmosphere:layers}{termosfer}.
\end{solution}

\begin{solution}{exo:g8:the-atmosphere:2}
(a) \qty{15}{\degreeCelsius}. (b) Kira-kira \qty{-56}{\degreeCelsius}.
(c) Kira-kira \qty{-3}{\degreeCelsius}. (d) Kira-kira
\qty{-86}{\degreeCelsius}. (e) Dari \qty{11}{km} sampai kira-kira
\qty{20}{km} (dan sepanjang beberapa kilometer di dekat \qty{50}{km} serta
tepat di atas \qty{86}{km}).
\end{solution}

\begin{solution}{exo:g8:the-atmosphere:3}
(a) Salah: suhu turun di \omterm{def:g8:the-atmosphere:layers}{troposfer}, naik di \omterm{def:g8:the-atmosphere:layers}{stratosfer}, turun di \omterm{def:g8:the-atmosphere:layers}{mesosfer},
dan naik lagi di \omterm{def:g8:the-atmosphere:layers}{termosfer}. (b) Benar. (c) Salah: \omterm{def:g8:the-atmosphere:ozone-layer}{lapisan ozon} terdapat di
\omterm{def:g8:the-atmosphere:layers}{stratosfer}. (d) Benar.
\end{solution}

\begin{solution}{exo:g8:the-atmosphere:4}
(a) \omterm{def:g8:the-atmosphere:layers}{Troposfer}. (b) \omterm{def:g8:the-atmosphere:layers}{Stratosfer} (dalam atmosfer baku, \omterm{def:g8:the-atmosphere:layers}{tropopause} berada pada
\qty{11}{km}). (c) \omterm{def:g8:the-atmosphere:layers}{Stratosfer}. (d) \omterm{def:g8:the-atmosphere:layers}{Mesosfer}.
\end{solution}

\begin{solution}{exo:g8:the-atmosphere:5}
(a) $20 - 6.5 \times 2 = 20 - 13 = \qty{7}{\degreeCelsius}$. (b) $20 -
6.5 \times 3 = 20 - 19.5 = \qty{0.5}{\degreeCelsius}$. (c) Udara harus
mendingin 20 derajat: $20 \div 6.5 \approx 3.1$, jadi pada kira-kira
\qty{3.1}{km}.
\end{solution}

\begin{solution}{exo:g8:the-atmosphere:6}
(a) $3 - 0.5 = \qty{2.5}{km}$. (b) $6.5 \times 2.5 = 16.25$: kira-kira 16
derajat. (c) $18 - 16.25 \approx \qty{2}{\degreeCelsius}$.
\end{solution}

\begin{solution}{exo:g8:the-atmosphere:7}
(a) \qty{500}{hPa}, \qty{250}{hPa}, dan \qty{125}{hPa}. (b) Di dekat tanah
aturan itu baik (\qty{500}{hPa} dibandingkan dengan \qty{506}{hPa}). Makin
tinggi, aturan itu memberikan nilai yang terlalu besar: \qty{250}{hPa}
dibandingkan dengan \qty{227}{hPa}, dan \qty{125}{hPa} dibandingkan dengan
\qty{96}{hPa}. Aturan separuh itu hanyalah pedoman kasar.
\end{solution}

\begin{solution}{exo:g8:the-atmosphere:8}
(a) $10~\text{hPa} = 10^1$~hPa; $10^3 \div 10^1 = 10^2$: seratus kali
lebih kecil. (b) $0.01~\text{hPa} = 10^{-2}$~hPa. (c) $10^3 \div 10^{-2} =
10^5$: seratus ribu kali lebih kecil.
\end{solution}

\begin{solution}{exo:g8:the-atmosphere:9}
(a) $6 \div 1.5 = 4$ kali lebih lebar. (b) Selagi balon itu naik, udara di
sekelilingnya makin lama makin lemah menekan, sehingga gas di dalamnya
mengembang dan balon itu membesar sampai kulitnya pecah. (c) Sesudah balon
meletus, parasut memperlambat jatuhnya radiosonde, sehingga radiosonde itu
tidak melukai siapa pun atau merusak apa pun.
\end{solution}

\begin{solution}{exo:g8:the-atmosphere:10}
(a) $20 \div 7 \approx 2.9$ kali lebih tebal. (b) Lebih dekat ke \omterm{def:g3:continents-and-oceans:poles}{kutub}:
\qty{11}{km} terletak \qty{4}{km} di atas \qty{7}{km} tetapi \qty{9}{km} di bawah \qty{20}{km} (atmosfer baku menggambarkan \omterm{def:g6:latitude-longitude:latitude}{lintang} tengah, bukan \omterm{def:g3:continents-and-oceans:equator}{khatulistiwa}).
(c) Pada awal penerbangan, \omterm{def:g8:the-atmosphere:layers}{tropopause} berada pada kira-kira \qty{7}{km}, di
bawah pesawat: pesawat itu berada di \omterm{def:g8:the-atmosphere:layers}{stratosfer}. Pada akhirnya, \omterm{def:g8:the-atmosphere:layers}{tropopause}
berada pada kira-kira \qty{20}{km}: pesawat itu berada di \omterm{def:g8:the-atmosphere:layers}{troposfer}.
\end{solution}

\begin{solution}{exo:g8:the-atmosphere:11}
(a) $1 - 506 \div 1013 \approx 0.50$: \textbf{50~\%}. (b) $1 - 227 \div
1013 \approx 0.78$: \textbf{78~\%}. (c) $1 - 12 \div 1013 \approx 0.99$:
\textbf{99~\%}. (d) Kira-kira 1~\%.
\end{solution}

\begin{solution}{exo:g8:the-atmosphere:12}
(a) $6371 \div 15 \approx 425$~km untuk setiap sentimeter. (b) $100 \div 425
\approx 0.24$~cm, kira-kira \qty{2.4}{mm}. (c) $11 \div 425 \approx
0.026$~cm, kira-kira \qty{0.26}{mm}. (d) $0.26 \div 0.1 \approx 2.6$: dua
atau tiga lembar kertas.
\end{solution}

\begin{solution}{pb:g8:the-atmosphere:1}
\textbf{1.} $300 \times 10 = 3000$~m.

\textbf{2.} $6000 \div 300 = 20$ menit; $\num{30000} \div 300 = 100$
menit.

\textbf{3.} 100 menit, yaitu 1~jam 40~menit.

\textbf{4.} Pada \qty{8}{km}: \omterm{def:g8:the-atmosphere:layers}{troposfer} (suhunya masih turun). Pada
\qty{30}{km}: \omterm{def:g8:the-atmosphere:layers}{stratosfer}.

\textbf{5.} $15 - (-50) = 65$ derajat.

\textbf{6.} $65 \div 10 = \qty{6.5}{\degreeCelsius}$ per kilometer.

\textbf{7.} Dari \qty{12}{km} sampai \qty{20}{km} suhunya tetap
\qty{-56}{\degreeCelsius}. Pendinginan berhenti antara \qty{10}{km} dan
\qty{12}{km}: \omterm{def:g8:the-atmosphere:layers}{tropopause} berada pada kira-kira \qty{11}{km}.

\textbf{8.} Suhunya naik: \qty{-55}{\degreeCelsius} pada \qty{22}{km},
\qty{-47}{\degreeCelsius} pada \qty{30}{km}. Balon itu berada di
\omterm{def:g8:the-atmosphere:layers}{stratosfer}, yang dihangatkan oleh ozon yang menyerap sinar ultraviolet
Matahari.

\textbf{9.} $1013 \div 2 \approx \qty{507}{hPa}$; tekanan itu tercapai
antara \qty{4}{km} (\qty{617}{hPa}) dan \qty{6}{km} (\qty{472}{hPa}).

\textbf{10.} $1013 \div 4 \approx \qty{253}{hPa}$.

\textbf{11.} Antara \qty{10}{km} (\qty{265}{hPa}) dan \qty{12}{km}
(\qty{194}{hPa}).

\textbf{12.} Dari \qty{10}{km}, tekanan harus turun $265 - 253 =
12$~hPa lagi, dari 71~hPa yang hilang sepanjang 2~km: $12 \div 71 \times 2
\approx 0.34$~km. Tekanan tinggal seperempat dari nilainya di \omterm{def:g4:world-physical-map:altitude}{permukaan
laut} pada kira-kira \textbf{\qty{10.3}{km}}.
\end{solution}
